\documentclass[10pt,a4paper]{amsart}
\usepackage[margin=19mm]{geometry}
\usepackage{amsmath,amssymb,mathtools}
\usepackage[scr=rsfs,cal=euler]{mathalfa}
\usepackage{xfrac}
\usepackage[unicode,hidelinks,bookmarks=false]{hyperref}
\newcommand{\ie}{i.e.\ }
\newcommand{\cf}{cf.\ }
\newcommand{\resp}{respectively\ }
\newcommand{\Subsection}{\subsection}
\providecommand{\nobreakdash}{\nobreak\hbox{-}}
\setlength{\emergencystretch}{2em}
\multlinegap0pt
\usepackage{bm}
\usepackage{bbm}
\usepackage{dsfont}
\usepackage{xspace}
\usepackage{cancel}
\usepackage{mathtools}
\usepackage{amsthm}
\makeatletter
\@ifundefined{theorem}{\newtheorem{theorem}{Theorem}}{}
\@ifundefined{claim}{\newtheorem{claim}[theorem]{Claim}}{}
\makeatother
\newcommand{\Prob}{\operatorname{Pr}}
\title[Maker playing against an invisible Breaker]{Maker playing against an invisible Breaker}
\author{Dennis Clemens, Fabian Hamann, Mirjana Mikala\v{c}ki, Yannick Mogge, Milo\v{s} Stojakovi\'c}
\date{Source: arXiv:2507.22519v1 (CC BY 4.0)}
\begin{document}
\maketitle

\noindent\emph{Source and licence.} Maker playing against an invisible Breaker. Version arXiv:2507.22519v1 (CC BY 4.0), distributed under the Creative Commons Attribution 4.0 International licence, as recorded in the arXiv licence metadata for this exact version and shown on the abstract page https://arxiv.org/abs/2507.22519v1. Full rights notice: licenses/arxiv-2507.22519-CC-BY-4.0.txt.\par\smallskip

\noindent\textit{Editorial scope.} Counted unit: the Claim whose source label is \texttt{claim:Ai1} in the article (source0.tex lines 1278--1282, source label \texttt{claim:Ai1}), delivered together with the article's own complete original proof of it (source0.tex lines 1284--1355, one \texttt{proof} environment of 72 lines). No mathematics is written, repaired or re-derived: statement and proof are the article's own text line for line. Required context retained uncounted: clm:connectivity\_number.rounds (lines 1193--1199). Every definition, display and label of those lines is retained unchanged. Omitted from this package and not redistributed: 1--1192, 1200--1277, 1283--1283, 1356--1942. Nothing inside a retained excerpt is omitted or altered: every retained body is delivered exactly as the article's lines are written, so no omission receipt of any kind accompanies this package. Citations: the retained excerpts carry no citation command at all, so no bibliography entry of the article is transcribed and no reference is redistributed. Reference adaptation: none was needed and none was made; every reference inside the delivered unit resolves inside it, so no unresolved reference or citation is produced. Printed slips kept literal: no printed word, symbol, hypothesis or number of the counted statement or of its proof was altered. Layout and class: the article's own class is \texttt{article} with packages including graphicx, hyperref, authblk, amsmath, amssymb, amsthm, mathrsfs, mathabx, dsfont, xcolor, enumitem, fontenc, lmodern, microtype; the wrapper is the lane's pinned \texttt{amsart} editorial wrapper at 10pt on A4, which restates the theorem-like environments the retained text needs and adds the article's own macro definitions that the retained text needs (\texttt{\textbackslash Prob}), with extra wrapper packages (bm, bbm, dsfont, xspace, cancel, mathtools). The delivered numbering is the unit's own. Assets: no retained span contains a figure environment, a graphics-inclusion command, a PDF-page inclusion, a table or a TikZ picture; no image file is copied into this package, the package declares no asset dependency and no image file is redistributed with it. Licence: version arXiv:2507.22519v1 of this article is distributed under the Creative Commons Attribution 4.0 International licence, as recorded in the arXiv licence metadata for this exact version and shown on the abstract page https://arxiv.org/abs/2507.22519v1. A full rights notice is delivered at licenses/arxiv-2507.22519-CC-BY-4.0.txt. Characters: every retained excerpt is pure ASCII, so the delivered engine, XeLaTeX with its default Latin Modern text font, reports no missing character.
\begin{claim}\label{clm:connectivity_number.rounds}
For $1 \leq i\leq \frac{1}{\varepsilon}$
the following holds:
$$
\Prob\left(\text{Maker does not forfeit in Stage $\mathcal{S}_i$, and $\mathcal{S}_i$ lasts more than $1.01\cdot \frac{\varepsilon n}{a}$ rounds} \right) = o(1).
$$
\end{claim}

\begin{claim}\label{claim:Ai1}
For every $i\leq \frac{1}{\varepsilon}$ we have
$\Prob(\overline{\mathcal{A}_{i,1}}|\mathcal{E}_{i-1})
=o(1)$.
\end{claim}

\begin{proof}
Assume that $\mathcal{E}_{i-1}$ holds.
By Claim~\ref{clm:connectivity_number.rounds}
we know that
with probability $1-o(1)$,
if Maker forfeits in Stage~$\mathcal{S}_{i}$,
then this must happen during the first 
$1.01\cdot \frac{\varepsilon n}{a}$ rounds of this stage. 

Now, let $Seq_j$ be any fixed sequence within the first 
$1.01\cdot \frac{\varepsilon n}{a}$ rounds
of Stage~$\mathcal{S}_{i}$. 
In order to succeed with a union bound over all
sequences in this stage, it is enough to show
that the probability of forfeiting during $Seq_j$ because of step (2) is upper bounded by $\exp(-n^{-0.1})$. Assume that $v$ is the vertex chosen in step (1) of sequence $Seq_j$, and let $C_v$ be its component. 

We first consider the case that $|C_v|\leq \frac{5b}{a}$. Then
let $C_1,\ldots,C_{\ell}$ be the components of $M[C_v]$
at the end of Stage $\mathcal{S}_{i-1}$.
Under assumption of $\mathcal{E}_{i-1}$
we have that $\bigwedge_{j\leq i-1} \mathcal{A}_{j,3}$
holds. Therefore, at the end of Stage $\mathcal{S}_{i-1}$ no such component
is dangerous, giving
$$
\sum_{u\in C_v} d_B(u)
=
\sum_{k=1}^{\ell} \sum_{u\in C_k} d_B(u)
<
\sum_{k=1}^{\ell} (|C_k|n - 0.3n)
\leq
|C_v|n - 0.3n.
$$ 
Since we now consider only the first $1.01\cdot \frac{\varepsilon n}{a}$ rounds of $\mathcal{S}_{i}$,
in the meantime Breaker can play at most
$1.01\cdot \frac{\varepsilon nb}{a}$ further edges, and hence increase the above sum of degrees by at most
$1.01\cdot \frac{\varepsilon nb}{a} + \binom{|C_v|}{2}<0.15n$, leading to
$
\sum_{u\in C_v} d_B(u)
<
|C_v|n - 0.15n
$
throughout sequence $Seq_j$.
Hence, more than $0.1n$ edges in $E(C_v,V\setminus C_v)$ are free throughout this sequence, and it follows
that the probability of forfeiting in step (2) is bounded from above by
\begin{align*}
\left( 1 - \frac{0.1n}{|C_v|\cdot (n-|C_v|)} \right)^{n^{0.2}}
& \leq 
\left( 1 - \frac{0.1n}{5bn/a} \right)^{n^{0.2}} 
=
\left( 1 - \frac{a}{50b} \right)^{n^{0.2}} 
<
\exp\left( - n^{0.1}\right)
\end{align*}
for large enough $n$, since $a$ and $b$ are constants.

If instead $|C_v|\geq n - \frac{5b}{a}$,
then we can calculate similarly,
as each component in $V\setminus C_v$ was not dangerous
at the end of Stage $\mathcal{S}_i$,
again leading to more than $0.1n$ free edges in $E(C_v,V\setminus C_v)$.

Hence, it remains to consider the case 
$\frac{5b}{a} < |C_v| < n-\frac{5b}{a}$.
Then $e(C_v,V\setminus C_v) > \frac{4bn}{a}$
for large enough $n$,
while we also have that 
$e_B(C_v,V\setminus C_v) < \frac{2bn}{a}$ during the considered sequence, since the game lasts at most $1.1\cdot \frac{n}{a}$ rounds in total. Hence, in this case
the probability of forfeiting in step (2) is bounded from above by
$0.5^{n^{0.2}} <
\exp\left( - n^{0.1}\right)$
for large enough $n$.
\end{proof}

\end{document}
